\section{Подпространства. Теоремы о размерности подпространства. Построение системы линейных алгебраических уравнений, ФСР которой является базисом подпространства}

\begin{definition}
Непустое подмножество $L$ линейного пространства $V$ называется \textbf{подпространством}, если оно само является линейным пространством относительно операций сложения и умножения на число, определенных в $V$. 
Критерий подпространства: $L$ --- подпространство тогда и только тогда, когда:
\begin{enumerate}
    \item $0 \in L$;
    \item $\forall x, y \in L$ и $\forall \lambda, \mu \in F$ выполняется $\lambda x + \mu y \in L$ (замкнутость относительно линейных комбинаций).
\end{enumerate}
\end{definition}

\begin{theorem}[О размерности подпространства]
Пусть $L$ --- подпространство конечномерного пространства $V$. Тогда:
\begin{enumerate}
    \item $\dim L \le \dim V$.
    \item Если $\dim L = \dim V = n$, то $L = V$.
\end{enumerate}
\end{theorem}

\begin{proof}
\begin{enumerate}
    \item Пусть $\dim L = k$. Возьмем базис $f_1, \dots, f_k$ в $L$. Так как $L \subset V$, эти векторы принадлежат $V$ и остаются линейно независимыми в $V$. Максимальное число линейно независимых векторов в $V$ равно $\dim V = n$. Следовательно, $k \le n$, то есть $\dim L \le \dim V$.
    \item Если $k = n$, то система из $n$ линейно независимых векторов $f_1, \dots, f_n$ в $n$-мерном пространстве $V$ образует базис всего пространства $V$. Любой вектор $x \in V$ линейно выражается через этот базис, а значит, принадлежит $L$. Следовательно, $V \subset L$, и так как $L \subset V$, получаем $L = V$.
\end{enumerate}
\hfill$\blacksquare$
\end{proof}

\begin{theorem}[О построении СЛАУ по базису подпространства]
Пусть подпространство $L \subset F^n$ задано своим базисом, состоящим из $k$ векторов $f_1, f_2, \dots, f_k$. Тогда можно построить однородную систему линейных алгебраических уравнений (СЛАУ) такую, что:
\begin{enumerate}
    \item Множество решений этой системы в точности совпадает с $L$.
    \item Исходные векторы $f_1, \dots, f_k$ образуют фундаментальную систему решений (ФСР) этой системы.
\end{enumerate}
\end{theorem}

\begin{proof}
Вектор $x = (x_1, x_2, \dots, x_n)^T \in F^n$ принадлежит подпространству $L$ тогда и только тогда, когда он линейно выражается через базисные векторы $f_1, \dots, f_k$. Это равносильно тому, что система векторов $\{f_1, \dots, f_k, x\}$ линейно зависима, или, что то же самое, ранг матрицы, составленной из этих векторов как столбцов, равен $k$ (рангу матрицы, составленной только из $f_1, \dots, f_k$).

Рассмотрим расширенную матрицу, где первые $k$ столбцов --- это координаты базисных векторов, а последний столбец --- координаты произвольного вектора $x$:
\[
B = \left( \begin{array}{cccc|c}
f_{11} & f_{12} & \dots & f_{1k} & x_1 \\
f_{21} & f_{22} & \dots & f_{2k} & x_2 \\
\vdots & \vdots & \ddots & \vdots & \vdots \\
f_{n1} & f_{n2} & \dots & f_{nk} & x_n
\end{array} \right)
\]
Применим к этой матрице элементарные преобразования строк, чтобы привести левый блок (размера $n \times k$) к ступенчатому виду. Так как ранг этого блока равен $k$, после приведения он будет иметь ровно $k$ ненулевых строк, а нижние $n - k$ строк будут состоять из одних нулей.

Для того чтобы вектор $x$ принадлежал $L$, ранг всей матрицы $B$ также должен оставаться равным $k$. Это возможно тогда и только тогда, когда в преобразованной матрице последние $n - k$ элементов последнего столбца (которые являются линейными комбинациями исходных $x_i$) также обращаются в ноль. 

Эти $n - k$ условий представляют собой однородную систему линейных уравнений относительно переменных $x_1, \dots, x_n$:
\[
\begin{cases}
c_{11}x_1 + c_{12}x_2 + \dots + c_{1n}x_n = 0 \\
\vdots \\
c_{n-k, 1}x_1 + c_{n-k, 2}x_2 + \dots + c_{n-k, n}x_n = 0
\end{cases}
\]
Матрица коэффициентов $C$ этой системы имеет размер $(n-k) \times n$ и ранг $n-k$ (так как преобразования строк не меняют линейную независимость уравнений). По теореме о структуре решений, размерность пространства решений этой системы равна $n - \operatorname{rang} C = n - (n - k) = k$. 

Поскольку все $k$ базисных векторов $f_1, \dots, f_k$ удовлетворяют этой системе (по построению) и они линейно независимы, они и образуют ФСР данной однородной системы.
\hfill$\blacksquare$
\end{proof}
